\title{LongRoPE: Extending LLM Context Window Beyond 2 Million Tokens}

\begin{abstract}
A substantial context window is a highly sought-after characteristic in large language models (LLMs). Nonetheless, the expansion of context length is currently constrained to approximately 128k tokens. This limitation stems from several challenges: the significant expense associated with fine-tuning, the limited availability of extensive textual datasets, and problematic token positions that can lead to catastrophic values.

In this work, we present LongRoPE, which for the first time expands the context window of pre-trained LLMs to an impressive \textbf{2048k} tokens, utilizing as few as 1k fine-tuning steps within training lengths up to 256k, while preserving the model's original performance on short contexts. This advancement is supported by three major contributions: (i) we discover and leverage two distinct non-uniformities in positional interpolation via an efficient search mechanism, resulting in superior initialization for fine-tuning and facilitating an 8$\times$ extension even without fine-tuning; (ii) we propose a staged extension approach, beginning with fine-tuning at 256k length, followed by a secondary positional interpolation on the extended LLM to reach the 2048k context window; (iii) we further recalibrate LongRoPE at an 8k length to restore its capability on short context tasks. Extensive evaluations conducted on LLaMA2 and Mistral over multiple benchmarks verify the efficacy of our technique. LongRoPE-extended models retain their original model architecture with only minor adjustments to the positional embedding, permitting the reuse of most existing optimizations. The code will be released at \texttt{https://github.com/microsoft/LongRoPE}

\section{Introduction}

Although Large Language Models (LLMs) have achieved notable advancements across a range of applications~\cite{gpt4,llama2}, they are frequently constrained by a fixed context window, such as the 4096-token ceiling in LLaMA2~\cite{llama2}. When input length surpasses this window, model accuracy diminishes, primarily because the model has not encountered these out-of-window positions during training. This limitation becomes particularly problematic for tasks that require processing many in-context examples~\cite{Huang2023FewerIM} or for the functioning of LLM agents~\cite{park2023generative,madaan2023selfrefine}.

Recent studies indicate that it is possible to increase a pre-trained LLM's context window up to approximately 128k tokens by fine-tuning the model on lengthier textual data~\cite{longlora,pi,yarn,zhang2024soaring,liu2023scaling}. However, extending the context window beyond this size is hindered by three main challenges. First, incorporating previously unseen position indices during extension introduces numerous abnormal values, which results in out-of-distribution complications and renders fine-tuning convergence problematic~\cite{pi}. This issue is exacerbated when expanding from 4k to over 1000k positions, since over 90\% of these positions are novel. Second, effective fine-tuning generally necessitates training data that matches the new, extended lengths. The availability of such lengthy texts—particularly those surpassing 1000k tokens—in existing datasets is minimal. In addition, the computational demands of processing such extensive sequences are considerable, often requiring excessive training durations and substantial GPU allocations. Third, with extremely large context window sizes, the model’s attention mechanism tends to distribute itself too broadly over the vast number of token positions. This dispersion results in diminished performance, especially when handling the shorter contexts for which the models were initially optimized~\cite{pi}.

To address the initial challenge, one strategy involves interpolating the RoPE positional embedding~\cite{rope,pi}, which adjusts novel position indices to fit within the pretraining positional spectrum, as illustrated in Fig.~\ref{fig:overview}. Position Interpolation (PI)~\cite{pi} modifies the rotary angles of RoPE by applying a linear interpolation determined by the extension ratio. The NTK method~\cite{ntk,dynamicntk} suggests using uneven interpolation and extrapolation schemes along the dimensions of RoPE. In contrast, YaRN~\cite{yarn} divides RoPE dimensions into three groups according to their frequencies and assigns to each either extrapolation, NTK, or linear interpolation, accordingly. Nevertheless, within the Transformer framework, positional embeddings carry a \emph{complex, spatially varying information entropy}. The nuanced patterns of non-uniformity are generally overlooked by current techniques, resulting in some information degradation, which ultimately restricts the achievable context window length.

Section~\ref{sec:analysis} demonstrates two principal empirical observations: \textbf{(1)} Optimal positional interpolation must address two distinct types of non-uniformities—namely, the variability across RoPE dimensions and across token positions. Specifically, dimensions with lower RoPE indices and early token positions require less interpolation, although the ideal approach is influenced by the desired extent of length extension. \textbf{(2)} By integrating these forms of non-uniformity into the interpolation process, one can better preserve the information encoded by the original RoPE, especially across critical dimensions and at key token positions. This approach reduces the informational degradation associated with positional interpolation, yielding superior initialization for downstream fine-tuning. In addition, it enables up to an 8$\times$ increase in context length, even without subsequent fine-tuning.

Building on these insights, we present \textbf{LongRoPE}, a robust approach capable of expanding the LLM context window to exceed 2 \emph{million} tokens. The design of LongRoPE centers around three principal advancements. First, {LongRoPE} leverages multidimensional non-uniformity within positional interpolation to its fullest extent. Specifically, it determines optimal rescale factors for the rotation angles in each RoPE dimension, taking into account the positions of tokens. Because the process of determining these rescale factors becomes exponentially more complex as the intended extension ratio increases, {LongRoPE} employs an evolutionary search algorithm enhanced by two optimization strategies to improve the efficiency of the search. An illustration of the resulting rescaled RoPE, obtained through this procedure, is provided in Fig.~\ref{fig:overview}.

Subsequently, {LongRoPE} employs a progressive and efficient extension methodology to reach a 2048k context window, circumventing the necessity for direct fine-tuning on extremely lengthy texts, which are both uncommon and difficult to obtain. Initially, the process identifies a context length of 256k on the pre-trained LLM and carries out fine-tuning at this length. Following this, since our non-uniform positional interpolation enables up to an 8$\times$ context length increase without requiring additional fine-tuning, a further search for novel RoPE rescaling factors is performed on the already fine-tuned and extended LLM. This procedure ultimately results in attaining a 2048k context window for LLaMA2 and Mistral~\cite{mistral}.

In order to prevent performance drops when handling the original, shorter context windows, {LongRoPE} maintains the adaptation of RoPE rescaling factors on the expanded LLM. Analogous to the approach used when scaling up context length from 256k to 2048k, we also reduce context windows to 4k and 8k on the LLM fine-tuned at 256k, utilizing our search algorithm to minimize positional interpolation. At inference time, if the sequence length falls below 8k, RoPE is revised with the optimized rescale factors obtained from the search process.

Comprehensive evaluations involving multiple LLMs and a range of long-context tasks confirm the efficacy of our approach. Our findings indicate that LongRoPE consistently achieves low perplexity across evaluation lengths spanning from 4k to 2048k, attains passkey retrieval accuracy exceeding 90\%, and matches the accuracy of established baselines on standard benchmarks constructed within a 4096-token context window. Furthermore, LongRoPE is broadly compatible with all LLMs utilizing RoPE embedding. We intend to make both our code and the LongRoPE-2048k models publicly available.

\section{Non-uniformity in Positional Interpolation}
\label{sec:analysis}

\subsection{Preliminary}

Transformer architectures necessitate the inclusion of explicit positional cues—typically through position embeddings—to encode the sequence of input tokens. In this study, we investigate RoPE~\cite{rope} position embedding, a prevalent technique adopted by contemporary LLMs. For a token located at position index $n$, the associated RoPE encoding can be concisely expressed as:

\begin{equation}
		\fontsize{7.8}{7.8} \selectfont
	\label{eq:rope}
	\left[ cos(n\theta_0), sin(n\theta_0), cos(n\theta_1), \cdot\cdot\cdot, cos(n\theta_{d/2-1}), sin(n\theta_{d/2-1}) \right]   
\end{equation}

where $d$ denotes the embedding dimension,
$n\theta_i$ specifies the rotary angle corresponding to the token at position $n$, and
$\theta_i = \theta^{-2i/d}$ encodes the rotation frequencies. RoPE typically employs a default value of 10000 for the base $\theta$.

\textbf{\emph{Context window extension ratio $s$} and positional interpolation}. We introduce $s$ as the quotient of the extended context length $L'$ and the initial context length $L$, given by $s=\frac{L'}{L}$.

In order to expand the context window from $L$ to $L'$, prevailing positional interpolation techniques advocate reducing the rotation frequencies $\theta_i$ according to the extension ratio $s$. Defining $\beta = \theta^{2/d}$ and introducing $\lambda$ as the effective rescale factor associated with $s$, we present a unified framework for these positional interpolation approaches as:

\begin{equation}	
	\fontsize{7.5}{7.5} \selectfont
	\label{eq:unified_rope}
	\begin{aligned}
		\left[cos\left(\frac{n}{\lambda(\beta)^0}\right), sin \left(\frac{n}{\lambda(\beta)^0}\right), cos\left(\frac{n}{\lambda(\beta)^1}\right), \cdot\cdot\cdot,  sin \left(\frac{n}{\lambda(\beta)^{d/2-1}}\right) \right]   
	\end{aligned}
\end{equation}

\textbf{Linear positional interpolation (PI)}. The PI approach~\cite{pi} involves applying a linear interpolation to position indices, constrained within the maximum length seen during pre-training. When expanding to a longer sequence with an extension ratio $s$, every positional rotation angle across RoPE dimensions is uniformly scaled down by a factor of $\lambda=s$. This process compresses the positional encoding, resulting in more densely packed position information and consequently making it harder for the model to differentiate tokens that are close together. As a result, the performance of PI typically degrades when used at larger extension ratios.

\textbf{NTK-based interpolation and extrapolation}. ~\cite{ntk,dynamicntk} approach RoPE from the perspective of information representation and utilize Neural Tangent Kernel (NTK) theory~\cite{jacot2018neural,tancik2020fourier}. To address the issue of position crowding in PI, they propose to redistribute interpolation effects among the RoPE dimensions, applying smaller scaling to dimensions associated with higher frequencies and larger scaling to those with lower frequencies. This leads to both positional interpolation and extrapolation, where $\lambda=s^{i}$. The enhanced dynamic NTK~\cite{dynamicntk} further modulates the extension ratio per position, taking the current sequence length into account. Unlike PI, which requires fine-tuning procedures, NTK-based techniques are capable of enlarging context windows without the need for fine-tuning, albeit typically constrained to a maximum extension ratio of 4$\times$.

\textbf{YaRN}~\cite{yarn} presents a notable enhancement in the effectiveness of positional interpolation. The method separates the RoPE dimensions into three distinct groups according to their frequency characteristics, assigning each group a specific interpolation approach. Dimensions corresponding to high frequencies are processed with extrapolation ($\lambda$=1), while those associated with low frequencies are interpolated linearly (PI). For intermediate frequencies, the NTK technique is utilized. The central innovation of YaRN resides in how it categorizes the RoPE dimensions; however, this grouping is determined through manual, empirical procedures. Consequently, this approach may yield less-than-optimal results when applied to newly developed LLMs.

\subsection{Study on Non-uniform Positional Interpolation}

Building upon the advancements introduced by NTK and YaRN, we observe that their performance improvements stem primarily from non-linear effects, particularly those that account for varying frequency patterns among RoPE dimensions to enable tailored approaches for interpolation and extrapolation. Nonetheless, prevailing methods for introducing non-linearity predominantly depend on manually crafted heuristics. This observation prompts two central inquiries: (1) Does the existing approach to positional interpolation represent the best possible solution? (2) Might there exist alternative forms of non-linearity that have yet to be investigated?

To investigate these issues, we employ evolution search (refer to Sec.~\ref{sec:method}) to identify improved non-uniform positional interpolation strategies for LLaMA2-7B. Perplexity serves as the optimization objective, evaluated on 5 randomly selected instances from the PG19~\cite{pg19} validation dataset. Our empirical study yields several notable insights.

\textbf{Finding 1}: \textit{RoPE dimensions exhibit substantial non-uniformities, which are not effectively handled by current positional interpolation methods.}

We optimize $\lambda$ individually for each RoPE dimension as described in Eq.~\ref{eq:unified_rope}. Table~\ref{tbl:exp1} presents a comparison of perplexity results for LLaMA2-7B, evaluated on the PG19 and Proof-pile~\cite{proof-pile} test sets, prior to any fine-tuning. The approach derived from our optimization yields notable gains, indicating that both linear (PI) and non-uniform (Dynamic-NTK and YaRN) interpolation techniques do not provide optimal results. Interestingly, on PG19, YaRN performs worse than both PI and NTK, failing to attain the intended context window length in models that have not been fine-tuned. For instance, with an 8k context length, YaRN’s perplexity increases substantially beyond 7k tokens.

In our investigation, the rescaling coefficients $\lambda$ in Eq.~\ref{eq:unified_rope} are non-uniform, in contrast to the constant scaling parameter $s$ utilized in PI, the analytic expressions in NTK, and the group-wise method in YaRN. The introduction of these varying rescaling factors leads to marked enhancements in LLaMA2's language modeling accuracy, measured by perplexity, for extended context lengths of 8k and 16k tokens, all achieved without additional fine-tuning. This benefit arises because the resultant positional encoding more faithfully maintains the integrity of the original RoPE structure—particularly in critical dimensions—which assists the LLM in more effectively discriminating between tokens in nearby positions.

\textbf{Finding 2}: \textit{RoPE for the initial tokens in the input sequence should be extrapolated with less interpolation.} 

We posit that for the first $\hat{n}$ tokens within the input sequences, the extent of interpolation applied by RoPE should be minimized. The rationale for this stems from the observation that these tokens tend to attract higher attention weights, thereby playing a pivotal role in the function of attention layers, as established in Streaming LLM~\cite{streamingllm} and LM-Infinite~\cite{han2023lminfinite}. To test this hypothesis, we experiment with context window extensions to 8k and 16k utilizing PI and NTK, and systematically refrain from interpolating the initial $\hat n$ tokens (where $\hat n$ = 0, 2, ..., 256). Notably, setting $\hat n$=0 corresponds to the standard PI and NTK without modification. Table~\ref{tbl:tokenposition} offers two main insights: \textbf{(1)} excluding position interpolation for the first tokens results in measurable performance gains; and \textbf{(2)} the ideal value for $\hat n$ is contingent on the particular context window length being targeted.

\textbf{Finding 3}: \textit{Non-uniform positional interpolation effectively extends LLM context window in both fine-tuning and non-fine-tuning settings.} 

Our research demonstrates that employing our optimized non-uniform position interpolation notably enhances model extension performance at 8k and 16k context lengths, even in the absence of fine-tuning. However, scaling to even longer context windows necessitates a fine-tuning step. Therefore, we further fine-tune LLaMA2-7B using the discovered RoPE scheme for a context window of 64k tokens (details in Appendix). As indicated in Table~\ref{tbl:finetune}, our approach yields substantial improvements over PI and YaRN methods, both prior to and following the fine-tuning of LLaMA2-7B. This advantage can be attributed to our precise implementation of non-uniform positional interpolation, which reduces information degradation and offers a superior starting point for the fine-tuning process.

\textbf{Summary}. In this work, we reveal two sources of non-uniformity: heterogeneity in RoPE dimensions and in token positions. Effectively harnessing these non-uniformities in the design of positional interpolation leads to marked advancements in large language model context extension.

\section{LongRoPE}

\label{sec:method}
Building on these insights, we propose LongRoPE. This method initially employs an effective search strategy to leverage both types of non-uniformities comprehensively, and subsequently applies this approach to expand the LLM context window to more than 2 million tokens.

\subsection{Problem Formulation}

These two sources of non-uniformity greatly expand the possible solution space, adding layers of difficulty to the optimization process. To manage these challenges, we recast the multidimensional non-uniform position interpolation task into a search-based optimization framework.

Consider a LLM designed for a context window of size $L'$ processing lengthy input sequences $\mathbf{X}$, where each element $\mathbf{x}\in\mathbf{X}$ contains more tokens than $L'$. For such inputs, the original rotary angle applied to the $i^{th}$ dimension within RoPE at token position $n$ is given by $\frac{n}{\beta^i}$. Under these definitions, the optimization problem can be posed as:

\begin{equation}
\begin{gathered}
		\label{eq:problem}

\underset {\mathbf{x}\in\mathbf{X}; \, |\mathbf{x}|\ge L'}{ \operatorname {arg\,min} } \,  \mathcal{L}\, \left( \text{LLM} (\text{RoPE}, \mathbf{X})\right),\quad \text{with}
\\
\scalebox{0.93}{
$\underset {\begin{subarray}{l} i=0,\cdot\cdot,\frac{d}{2}-1;\\ n\in[ 0,|\mathbf{x}|);\end{subarray}}{\text{RoPE}_{(n)}} = {\left[\cdots, \cos\left(\mathbb{I}(\hat{\lambda}_{i}, \hat{n})\frac{n}{\beta^i}\right),\, \sin\left(\mathbb{I}(\hat{\lambda}_{i}, \hat{n})\frac{n}{\beta^i}\right),\cdots\right] }$}\\
	\text{with} \quad \mathbb{I}(\hat{\lambda}_{i},\hat{n}) =
	\begin{cases}
	1&\mbox{\quad if $n<\hat{n}$}\\
	\frac{1}{\lambda_{i}} &\mbox{\quad if $n\geq\hat{n}$}
	\end{cases}
	\end{gathered}
\end{equation}

In this approach, we define a collection of rescale factors, $\mathbb{I}(\hat{\lambda}_{i},\hat{n})$, to address both types of non-uniformities. Here, $\hat{\lambda}_i$ represents irregularities across RoPE dimensions, while $\hat{n}$ corresponds to non-uniformity with respect to token positions. More precisely, $\mathbb{I}(\hat{\lambda}_{i},\hat{n})$ serves to adjust the rotation angle for the $i^{th}$ RoPE dimension, with $\hat{\lambda}_i$ acting as the rescale parameter and $\hat{n}$ indicating the position threshold. For the first $\hat{n}-1$ token positions, the scaling factor $\hat{\lambda}_i$ is not applied and the standard RoPE rotary angle $\frac{n}{\beta^i}$ is maintained. Once token positions reach $n\geq\hat{n}$, the rescaling factor comes into effect.

For a given target context window of size $L'$, our goal is to determine the set of optimal rescale factors, specifically ($\mathbb{I}(\hat{\lambda}_{0},\hat{n})$, $\mathbb{I}(\hat{\lambda}_{1},\hat{n})$, ... , $\mathbb{I}(\hat{\lambda}_{i},\hat{n})$, ...), corresponding to each RoPE dimension from the $1^{st}$ up to the $d^{th}$. Implementing these rescale factors in the target $\text{LLM}$, thereby modifying its $\text{RoPE}$, is intended to yield the lowest possible next-token prediction loss $\mathcal{L}$ (i.e., perplexity) when processing input data $\mathbf{X}$ composed of $L'$ tokens.

\subsection{Searching the Non-uniform Position Interpolation}

To address the challenge outlined in Eq.~\ref{eq:problem}, we present a straightforward yet powerful approach that identifies optimal RoPE rescaling parameters, thereby leveraging the complex, multidimensional disparities inherent in positional encodings.

\textbf{Search space}. Our method constructs an extensive search space to accommodate both forms of non-uniformity. This is detailed in Table~\ref{tbl:searchspace}. More precisely, the approach enables the independent optimization of rescale factors for each individual dimension within the RoPE embedding. For ease of formulation, we opt to search over $\lambda_i$ and $\hat{n}$, in lieu of directly exploring $\mathbb{I}(\hat{\lambda}_{i},\hat{n})$, where $\hat{\lambda}_{i}=1/\lambda_i$. As depicted in Table~\ref{tbl:searchspace}, the value of $\lambda_i$ can be selected from as low as 1.0 (representing simple extrapolation) up to $s\times1.25$ (which affords a degree of interpolation beyond that used in PI), using increments of 0.01, with $s$ denoting the targeted extension ratio of the context window.

$\hat{n}$ determines how many of the earliest token positions maintain the unaltered RoPE embedding, foregoing position interpolation. In practice, $\hat{n}$ is selected from the set \{0, 1, 2, 4, 8, 12, 16, 20, 24, 28, 32, 64, 128, 256\}. Setting $\hat{n}=0$ results in every token position applying the learned rescale factors.

\textbf{Evolution-based search}. The search space outlined in Table~\ref{tbl:searchspace} encompasses a vast number of positional interpolation configurations, creating substantial complexity for efficient navigation. As an illustration, an extension factor of $s=4\times$ results in $400^{128/2}\times14 = 4\times10^{167}$ possible combinations. The number of alternatives grows exponentially with higher extension ratios. To efficiently cope with this immense space, we employ evolution search~\cite{guo2020single} and implement two optimization strategies that significantly enhance the efficiency of the search process. The overall methodology is depicted in Algorithm~\ref{alg:search}.

\textit{Optimized initial population generation}. Rather than selecting all $P$ rescale factors at random to create the initial population, we explicitly include the three RoPE rescale factors associated with PI, NTK, and YaRN as members of this group. The other $P$-3 individuals are generated by randomly mutating these three rescale factors with probability $p$.

\textit{Monotonically non-decreasing constraint}. Once the initial population is established, we evaluate the LLM perplexity associated with each member. For this, we assign the designated RoPE rescale factors to the target LLM and assess perplexity on the input $\mathbf{X}$. Individuals ranking in the top-k are then chosen as parents for subsequent evolutionary steps. Due to the expansive search space, straightforward application of mutation and crossover operations may frequently result in suboptimal candidates, incurring a high number of redundant perplexity evaluations. This inefficiency becomes especially pronounced for large $L'$, since each perplexity assessment is computationally demanding.

To tackle this, we enforce a constraint of non-decreasing monotonicity on the RoPE rescaling factors being sampled, formally expressed as $\lambda_i\le \lambda_{i+1}$. Only those RoPE variants adhering to this condition are utilized for LLM perplexity evaluation, thereby greatly lowering the computational burden of the search. Concretely, the requirement stipulates that the sequence $\lambda_i$ should be monotonically increasing as a function of the RoPE dimension index ($i=0,\ldots,63$). This monotonicity across dimensions is motivated by the NTK theory~\cite{jacot2018neural,tancik2020fourier,ntk}, which implies that lower-dimensional, higher-frequency components demand less interpolation (hence, a smaller $\lambda_i$), whereas higher-dimensional, lower-frequency components permit more interpolation (corresponding to a larger $\lambda_i$).

\begin{algorithm}[t]
	\small
	\caption{The search algorithm}
	\textbf{Input:} target LLM, input samples $\mathbf{X}$, population size $P$, mutation size $N_1$, crossover size $N_2$, maximum iterations $\mathcal{T}$, mutation probability $p$\\

	\begin{algorithmic}[1]
		\label{alg:search}
		\STATE Top-k $\leftarrow \phi$;
		\STATE $\text{P}_0 \leftarrow$ \textit{Initialize\_population\_with\_optimization}($P$, $\mathbf{X}$, $p$);
		\FOR{$i$ from $1$ to $\mathcal{T}$}
		\STATE \textit{Compute\_perplexity} (LLM, $\text{P}_{i-1}$, $\mathbf{X}$);
		\STATE Top-k $\leftarrow$ \textit{Update\_Topk}(Top-k, $\text{P}_{i-1}$);
		\STATE $\text{P}_{mutation} \leftarrow$ \textit{Mutation\_with\_mono\_constraint}(Top-k, $N_1$, $p$);
		\STATE $\text{P}_{crossover} \leftarrow$ \textit{Crossover\_with\_mono\_constraint}(Top-k, $N_2$);
		\STATE $\text{P}_i \leftarrow \text{P}_{mutation} \cup \text{P}_{crossover} \cup \text{Top-k}$;
		\ENDFOR
		\STATE Return the member of Top-k with the minimum perplexity;
	\end{algorithmic}
\end{algorithm}

\textbf{8$\times$ extension without fine-tuning}. By leveraging evolutionary search, we successfully determine non-uniform RoPE rescale factors that safeguard essential dimensions and positions, thereby reducing the loss of information typically caused by interpolation. As shown in Fig.\ref{fig:non-ft}, our approach enables LLaMA2 to achieve an extended context window of 32k from its original 4k limit, all without the need for fine-tuning. Conversely, approaches like PI and the non-uniform versions of NTK and YaRN result in a notable increase in perplexity after only a 2$\times$ context extension.

\subsection{Extending LLM Context Window to 2048K}
\label{sec:extendto2048k}

\textbf{Progressive extension to 2048k}. In this section, we present our approach for expanding the context window of pre-trained LLMs from the standard 4k length up to beyond 2048k. Our findings show that employing non-uniform positional interpolation enables an 8$\times$ window increase without the need for fine-tuning. However, when aiming for substantially larger expansions (such as a 512$\times$ increase), fine-tuning becomes essential. One approach is to determine suitable RoPE rescaling factors specifically for the target 2048k window size prior to fine-tuning. Despite this, the approach is hindered by the substantial computational demands of training. Additionally, based on our observations, fine-tuning LLMs with such high extension ratios proves to be quite difficult (see Appendix for further details).

Fortunately, LongRoPE demonstrates efficacy with both base and fine-tuned extended LLMs. Consequently, we propose a resource-efficient, incremental approach that attains the desired 2048k context window using only 1k fine-tuning steps, all while restricting the training sequence length to 256k.

$\Diamond$ \textit{LongRoPE-based search for scaling pre-trained LLMs to 256k}. Using LLaMA2 for demonstration, we perform search targeting context windows of 128k and 256k. This results in extension factors of 32$\times$ and 64$\times$, respectively.

$\Diamond$ \textit{Fine-tuning to 256k}. Subsequently, we adapt the pre-trained LLM to support a 256k context window through fine-tuning. In particular, our approach involves initially fine-tuning LLaMA2 for 400 iterations utilizing RoPE rescaled factors set to 128k. After completing this phase, we switch the RoPE rescaled factors to 256k within the obtained checkpoint, followed by an additional 600 fine-tuning steps. This staged procedure demonstrates greater efficiency compared to attempting immediate fine-tuning at the 256k setting.

$\Diamond$ \textit{Scaling fine-tuned extended LLM to 2048k using LongRoPE search}. In the last stage, an additional search is applied to the LLM that has already been fine-tuned for 256k token sequences. This procedure enables a substantial enlargement of the context window, reaching up to 2048k tokens, all achieved without any additional fine-tuning steps. Consequently, the context window size is expanded by a factor of $512\times$.

\textbf{Recovery for short context windows}. Upon expanding the context window to an extensive length of 2048k, we observe a decline in performance within the initial context window range. This phenomenon is well-documented in the literature on positional interpolation~\cite{pi}, which compels the position embeddings for higher-index tokens within the initial context to be confined to a substantially limited segment of the embedding space, thereby impairing the language model's effectiveness. Notably, when applying a 512$\times$ scaling factor, the positions corresponding to the original 4k context window are compressed into an even more congested area.

To address this, we conduct an additional evolutionary search on the expanded LLM, fine-tuning the RoPE rescale factors specifically for shorter context lengths (such as 4k and 8k). For these reduced lengths, we limit the upper bound of the searched $\lambda$ values, since there is a decreased need for positional interpolation. At inference time, the LLM automatically updates the relevant RoPE rescale factors in response to the input context length.

\section{LongRoPE}

\label{sec:method}
Building upon these observations, we propose LongRoPE. This approach begins by implementing an efficient search algorithm designed to capitalize on the two observed non-uniformities. Leveraging this algorithm, LongRoPE successfully extends the LLM context window to support over 2 million tokens.

\subsection{Problem Formulation}

These two types of non-uniformity significantly expand the solution space and complicate the optimization process. To tackle this challenge, we recast the problem of optimizing multidimensional non-uniform position interpolation as a search problem.

Given an LLM with a target context window of size $L'$ and a set of lengthy input documents $\mathbf{X}$—where each sequence $\mathbf{x} \in \mathbf{X}$ contains more tokens than $L'$—the original rotary angle assigned to the $i^{th}$ dimension in the RoPE embedding for a token at position $n$ is given by $\frac{n}{\beta^i}$. Accordingly, we define the optimization problem as follows:

\begin{equation}
\begin{gathered}
		\label{eq:problem}

\underset {\mathbf{x}\in\mathbf{X}; \, |\mathbf{x}|\ge L'}{ \operatorname {arg\,min} } \,  \mathcal{L}\, \left( \text{LLM} (\text{RoPE}, \mathbf{X})\right),\quad\text{with}
\\
\scalebox{0.93}{
$\underset {\begin{subarray}{l} i=0,\cdot\cdot,\frac{d}{2}-1;\\ n\in[ 0,|\mathbf{x}|);\end{subarray}}{\text{RoPE}_(n)}= {\left[\cdots,\cos\left(\mathbb{I}(\hat{\lambda}_{i}, \hat{n})\times\frac{n}{\beta^i}\right), \sin\left(\mathbb{I}(\hat{\lambda}_{i}, \hat{n})\times\frac{n}{\beta^i}\right),\cdots\right] }$}\\
	\text{where}\quad \mathbb{I}(\hat{\lambda}_{i},\hat{n})=
	\begin{cases}
	1 &\text{if $n<\hat{n}$}\\
	{\frac{1}{\lambda}_{i}} &\text{if $n\geq\hat{n}$}
	\end{cases}
	\end{gathered}
\end{equation}

We define a collection of rescale factors, $\mathbb{I}(\hat{\lambda}_{i},\hat{n})$, to address two distinct sources of non-uniformity. Here, $\hat{\lambda}_i$ corresponds to dimension-specific non-uniformity in RoPE, while $\hat{n}$ captures the non-uniformity associated with token positions. In practice, the function $\mathbb{I}(\hat{\lambda}_{i},\hat{n})$ serves to modify the rotation angle for the $i^{th}$ dimension of RoPE: the parameter $\hat{\lambda}_i$ acts as the scaling coefficient, and $\hat{n}$ denotes the position threshold for applying the scaling. For the initial $\hat{n}-1$ token positions, we do not modify the original RoPE rotary angle $\frac{n}{\beta^i}$. Once the token position reaches $n \geq \hat{n}$, the rescaling factor begins to influence the rotation angle.

For a specified target context window length $L'$, our goal is to determine the best set of rescale factors ($\mathbb{I}(\hat{\lambda}_{0},\hat{n})$, $\mathbb{I}(\hat{\lambda}_{1},\hat{n})$, ... , $\mathbb{I}(\hat{\lambda}_{i},\hat{n})$, ...) across all RoPE dimensions from $1^{st}$ to $d^{th}$. By applying these rescaled factors to the target $\text{LLM}$’s $\text{RoPE}$, we aim to minimize the next token prediction loss $\mathcal{L}$ (representing perplexity) over input instances $\mathbf{X}$ with token sequence length $L'$.

\subsection{Searching the Non-uniform Position Interpolation}

To address the challenge posed in Eq.~\ref{eq:problem}, we propose an intuitive and effective approach that identifies the most suitable RoPE rescale factors, thereby harnessing the inherently multidimensional, non-uniform structure of positional embeddings to the fullest.

\textbf{Search space}. Our method constructs a comprehensive search space designed to capture both types of non-uniformity. This is outlined in Table~\ref{tbl:searchspace}. Concretely, we permit the optimization process to assign unique rescale factors to each dimension within the RoPE embedding. To streamline the search, our algorithm explores $\lambda_i$ and $\hat{n}$ directly, rather than targeting $\mathbb{I}(\hat{\lambda}_{i},\hat{n})$, with $\hat{\lambda}_{i}$ defined as $1/\lambda_i$. Table~\ref{tbl:searchspace} details that $\lambda_i$ can be selected from a range starting at 1.0 (representing straightforward extrapolation) and extending up to $s\times1.25$ (which allows for broader interpolation than PI), with increments of 0.01. Here, $s$ denotes the intended extension ratio for the context window.

The parameter $\hat{n}$ specifies how many leading token positions maintain the original RoPE embedding, foregoing position interpolation. In practice, $\hat{n}$ is selected from the set \{0, 1, 2, 4, 8, 12, 16, 20, 24, 28, 32, 64, 128, 256\}. If $\hat{n}=0$, then the learned rescale factors are applied to every token position.

\textbf{Evolution-based search}. The search space described in Table~\ref{tbl:searchspace} encompasses a vast array of positional interpolation options, making efficient search particularly difficult. For instance, an extension factor of $s=4\times$ results in $400^{128/2}\times14$=4$\times10^{167}$ potential configurations. As the extension ratio increases, the size of the search space grows exponentially. To cope with this complexity, we adopt evolution search~\cite{guo2020single} and incorporate two optimization strategies that significantly enhance the efficiency of the search process. The complete search workflow is outlined in Algorithm~\ref{alg:search}.

\textit{Optimized initial population generation}. Rather than relying solely on random initialization for the population of $P$ rescale factors, our approach explicitly includes the three RoPE rescale factors—PI, NTK, and YaRN—as distinct members in the initial population. The rest of the $P$-3 individuals are produced by randomly mutating these three rescale factors, each with a mutation probability of $p$.

\textit{Monotonically non-decreasing constraint}. Following the creation of the initial population, we evaluate the LLM perplexity of each candidate. This involves applying each individual’s specified RoPE rescale factors within the target LLM and then determining the perplexity of the input $\mathbf{X}$. The k candidates with the lowest perplexity scores are selected as parents for the evolutionary process. Nonetheless, due to the immense size of the search space, simple mutation and crossover may lead to exploration of suboptimal candidates, resulting in superfluous perplexity evaluations. This inefficiency becomes even more pronounced as $L'$ increases, since each perplexity inference is computationally expensive.

To tackle this issue, we enforce a monotonic non-decreasing constraint on the sampled RoPE rescaling factors such that $\lambda_i \le \lambda_{i+1}$. Only those RoPE rescaling configurations that fulfill this constraint are utilized during perplexity evaluation on LLMs, which effectively narrows the search space. In particular, we stipulate that $\lambda_i$ must increase monotonically with respect to the RoPE dimension index ($i=0,\ldots,63$). This requirement draws motivation from NTK theory~\cite{jacot2018neural,tancik2020fourier,ntk}, which indicates that lower-dimensional components—corresponding to higher frequencies—necessitate less interpolation (i.e., a smaller $\lambda_i$), whereas higher-dimensional components—representing lower frequencies—can accommodate more interpolation (i.e., a larger $\lambda_i$).

\begin{algorithm}[t]
	\small
	\caption{The search algorithm}
	\textbf{Input:} target LLM, input samples $\mathbf{X}$,   population size $P$, mutation size $N_1$, crossover size $N_2$, max iterations $\mathcal{T}$, mutate probability $p$\\

	\begin{algorithmic}[1]
		\label{alg:search}
		\STATE Initialize Top-k as empty;	
		\STATE $\text{P}_0$ = \textit{Initialize\_population\_with\_optimization} ($P$, $\mathbf{X}$, $p$); 
		\FOR{$i$ from $1$ to $\mathcal{T}$}
		\STATE \textit{Compute\_perplexity} (LLM, $\text{P}_{i-1}$, $\mathbf{X}$);
		\STATE  Top-k $\leftarrow$ \textit{Update\_Topk} (Top-k, $\text{P}_{i-1}$);
		\STATE $\text{P}_{mutation}$ $\leftarrow$ \textit{Mutation\_with\_mono\_constraint} (Top-k, $N_1$, $p$); 
		\STATE $\text{P}_{crossover}$ $\leftarrow$ \textit{Crossover\_with\_mono\_constraint} (Top-k, $N_2$);
		\STATE $\text{P}_i$ = $\text{P}_{mutation}$ $\cup$ $\text{P}_{crossover}$ $\cup$ Top-k;
		\ENDFOR
		\STATE Return the candidate in Top-k with minimal perplexity;
	\end{algorithmic}
\end{algorithm}

\textbf{8$\times$ extension without fine-tuning}. Through our evolutionary search approach, we successfully discover non-uniform RoPE rescale factors that selectively preserve important dimensions and positional encodings, thereby reducing the loss of information resulting from interpolation. As illustrated in Fig.\ref{fig:non-ft}, this strategy enables the context window of LLaMA2 to be scaled from 4k to 32k tokens without requiring fine-tuning. By comparison, other approaches like PI, as well as non-uniform NTK and YaRN, experience substantial increases in perplexity beyond a 2$\times$ window extension.

\subsection{Extending LLM Context Window to 2048K}
\label{sec:extendto2048k}

\textbf{Progressive extension to 2048k}. In this section, we present our approach for scaling the context window of pre-trained LLMs from the standard 4k length up to more than 2048k. Our experiments indicate that utilizing non-uniform positional interpolation enables up to 8$\times$ expansion of the context window without the need for fine-tuning. To achieve even greater extensions, such as a 512$\times$ increase, fine-tuning becomes indispensable. One strategy involves identifying suitable RoPE rescaling parameters specific to the desired 2048k context length, followed by fine-tuning. Nonetheless, this process is confronted by significant demands on computational resources for training. Additionally, our empirical observations suggest that it remains difficult to effectively fine-tune LLMs when faced with such substantial extension ratios (refer to the Appendix for details).

Luckily, LongRoPE proves successful when applied to both the base and further-trained extended LLM variants. Consequently, we present a streamlined, incremental approach that reaches the intended 2048k target using only 1k fine-tuning iterations, all conducted with training sequences limited to 256k in length.

$\Diamond$ \textit{Expanding the context length of pre-trained LLMs to 256k using LongRoPE search}. Using LLaMA2 as a case study, we perform a search targeting context window sizes of 128k and 256k. The corresponding extension factors achieved during this process are 32$\times$ and 64$\times$, respectively.

$\Diamond$ \textit{Fine-tuning to 256k}. Subsequently, we adapt the pre-trained LLM to reach a context window of 256k through further fine-tuning. To do so, we initially fine-tune LLaMA2 over 400 iterations employing RoPE rescaled factors corresponding to 128k. Upon completing these steps, we update the RoPE rescaled factors to target 256k on the resulting checkpoint and proceed with 600 more fine-tuning steps. This staged approach is more effective than immediate fine-tuning for 256k.

$\Diamond$ \textit{Scaling up a fine-tuned extended LLM to 2048k using LongRoPE search}. In the final stage, we apply an additional search process to the LLM that has already been fine-tuned with a 256k context length. This procedure enables us to achieve a significantly expanded context window of 2048k tokens, all without the need for additional fine-tuning steps. The overall expansion ratio attained is $512\times$.

\textbf{Restoration in shorter context windows}.  Upon expanding the context window to an extensive 2048k tokens, we observe a degradation in performance for sequences within the initial context window range. This phenomenon is well-documented in positional interpolation~\cite{pi}, where extending the range compresses the position embeddings of the original window into a more limited subspace, which detrimentally impacts model performance. Specifically, with a 512$\times$ extension, the position representations for the initial 4k tokens become densely packed.

To address this issue, we conduct an additional evolutionary search on the expanded LLM in order to fine-tune RoPE rescale factors tailored to shorter context lengths (such as 4k and 8k). For these shorter lengths, we constrain the upper limit of the searched $\lambda$ values, reflecting the reduced necessity for positional interpolation. At inference time, the LLM adaptively updates the associated RoPE rescale factors as needed.

 \section{Experiments}

\textbf{Evaluation Tasks and Models}. LongRoPE is implemented on both LLaMA2-7B and Mistral-7B models. The evaluation framework covers three main areas: (1) assessing the perplexity scores of long-context language models on extensive textual inputs; (2) conducting a Passkey retrieval task, which tests each model’s proficiency in identifying a designated passkey amidst a large quantity of distractor content; and (3) benchmarking performance on conventional LLM tasks using a short context window of 4096 tokens.

\textbf{Fine-tuning}. For the LLaMA2 model, we employ a learning rate of $2 \times 10^{-5}$ with a linear decay schedule and a global batch size set to 32. Fine-tuning is conducted for 400 steps on the Redpajama~\cite{redpajama} dataset, which is partitioned into segments of 128k tokens, each marked by BOS and EOS tokens. Subsequently, continuing from the obtained checkpoint, we further train for 600 steps to extend the context window to 256k. Training for the 128k context is carried out on 8 A100 GPUs utilizing the distributed training framework~\cite{cube}, whereas scaling up to 256k context necessitates 16 A100 GPUs. For Mistral, a fixed learning rate of $1 \times 10^{-6}$ and a global batch size of 64 are adopted. Consistent with the approach in YaRN~\cite{yarn}, both 128k and 256k context models are trained for 400 steps on Together Computer's Long-Data Collections~\cite{mistraldata}, with a sequence length of 16k. Training in this case makes use of 4 A100 GPUs.

\textbf{Search}. For target window sizes up to 256k, the following parameters are applied: $P$=64, $N_1$=$N_2$=16, $p$=0.3, and $\mathcal{T}$=40. During each iteration, the top-32 candidates are selected for mutation and crossover. Perplexity evaluation employs five randomly chosen samples from the PG19 validation set, each sample meeting the minimum requirement of the intended context length. For window sizes exceeding 512k, the sizes of the population, mutation, and crossover pools are reduced by half. In this case, perplexity is assessed on three randomly selected validation samples from Pile-Books3~\cite{pile}.

\textbf{Baselines}. To obtain a 2048k context length, we conducted fine-tuning of models with 128k and 256k window sizes, resulting in LongRoPE-2048k (ft=128k) for LLaMA2 and LongRoPE-2048k (ft=256k) for Mistral. Our analysis includes a comparison between these four models and leading context window extension baselines, namely publicly available LLMs that have undergone fine-tuning after positional interpolation through methods such as PI, NTK, and YaRN. The set of baselines encompasses Together-32k~\cite{together}, Code LLaMA~\cite{codellama}, LongLoRA-full-FT-100k~\cite{longlora}, YaRN-LLaMA, and YaRN-Mistral~\cite{yarn}.

\subsection{Main Results}

\label{sec:mainresult}
\textbf{Language modeling of long sequences up to 256k}. Our initial experiments focus on assessing performance relative to leading extended LLMs, constrained to a maximum sequence length of 256k during evaluation. To illustrate the robustness of our approach, we employ two benchmarks: the Proof-pile~\cite{pg19} and PG19~\cite{pile} test sets. Perplexity is measured across different context windows using a fixed sliding window of 256 tokens. The complete test partition, comprising 100 documents, is used for PG19. In alignment with YaRN~\cite{yarn}, we select 10 samples from Proof-pile at random, ensuring that each selected document contains at least 128k tokens.

Table~\ref{tbl:proof-pile} and Table~\ref{tbl:pg19} present a comparison of perplexity scores for LLaMA2 and Mistral models, each extended using various interpolation strategies, evaluated on the Proof-pile and PG19 datasets, respectively. Two main findings emerge from these results: \textbf{(1)} Our extended models consistently exhibit a reduction in perplexity as evaluation length increases from 4k to 256k, demonstrating enhanced capability in utilizing extended contexts. \textbf{(2)} Notably, the LongRoPE-2048k variants maintain superior performance compared to leading baseline models within a 256k context window, even when operating with context lengths up to 16$\times$ greater—an environment that typically hampers short-length performance.

\textbf{Modeling language on sequences surpassing 2000k tokens}. To assess performance on exceptionally lengthy texts, we utilize the Books3~\cite{pile} dataset. For practical evaluation, 20 books are randomly chosen, each containing more than 2048k tokens, and a sliding window approach of 256k is applied.

As indicated in Table~\ref{tbl:books3}, LongRoPE is effective in expanding the context window of both LLaMA2-7B and Mistral-7B models up to 2048k, while maintaining perplexity levels on par with, or superior to, the baseline methods for shorter contexts ranging from 8k to 128k. Notably, there are substantial differences in the performance of the 2048k variants of LLaMA2 and Mistral. Mistral demonstrates stronger performance than baseline models at lower context lengths, yet its perplexity rises above 7 when the context extends past 256k. LLaMA2 exhibits anticipated behavior, with perplexity generally improving for longer contexts, except for minor increases observed at 1024k and 2048k. Furthermore, for LLaMA2, LongRoPE-2048k yields better results when fine-tuned at 256k as opposed to 128k, attributable to a smaller secondary extension ratio (8$\times$ versus 16$\times$). Conversely, Mistral attains optimal outcomes with a fine-tuning window of 128k. This difference largely arises because, in the fine-tuning process for Mistral at 128k and 256k, we adopt a 16k training sequence length as recommended by YaRN, which limits Mistral’s capacity for extending the context window after fine-tuning.

\textbf{Passkey retrieval}. In this section, we investigate the practical context window length for generation tasks. We adopt the synthetic passkey retrieval benchmark introduced by~\cite{passkey}. In this setup, the model's objective is to locate a randomly embedded passkey—defined as a five-digit number—within an extended document. Details of the prompt format are provided in the appendix. We conduct the passkey retrieval experiment over 10 runs, each time positioning the passkey at a uniformly random point within the given context window for evaluation.

Fig.~\ref{fig:passkey} presents a comparison of retrieval accuracy with various baseline methods. The accuracy of existing models declines sharply to zero once the sequence length exceeds 128k. By contrast, even in the demanding scenario of extracting a passkey from sequences containing millions of tokens, our LongRoPE-LLaMA2-2048k (ft=256k) consistently preserves high retrieval accuracy ($\ge$90\%) across the entire range from 4k up to 2048k. LongRoPE-Mistral-2048k (ft=128k) maintains perfect accuracy up to 1800k tokens, before dropping to 60\% at 2048k, which is consistent with the trend noted in Table~\ref{tbl:books3}, where perplexity exhibits a slight increase at the longest context window.

\textbf{Evaluation on standard benchmarks within the initial context window}. The LongRoPE-2048k models are assessed on their original context length utilizing the Hugging Face Open LLM Leaderboard~\cite{huggingfaceboard} under both zero-shot and few-shot paradigms. Specifically, our evaluation protocol includes a 25-shot setting for ARC-Challenge~\cite{allenai:arc}, a 10-shot setup for HellaSwag~\cite{zellers2019hellaswag}, 5-shot MMLU~\cite{mmlu}, and a zero-shot configuration for TruthfulQA~\cite{lin2021truthfulqa}.

According to Table~\ref{tbl:commonsense}, our models demonstrate similar performance to the original benchmarks, which were intended for shorter context lengths, and actually surpass the baseline Mistral model on TruthfulQA by a margin of +0.5\%. Although LongRoPE-LLaMA2-2048k, which was fine-tuned with a 256k context, exhibits a marginally greater decline in performance, it generally maintains satisfactory results across the majority of evaluated tasks.

\subsection{Ablation Results}

\textbf{Effectiveness of the second positional interpolation}. Within the framework of our progressive extension strategy, we apply our search-based approach to perform a second round of non-uniform positional interpolation on large language models that have already undergone fine-tuning. To assess the impact of this method, we conduct experiments utilizing our fine-tuned LLaMA2-256k model as a baseline. This model is further scaled to context lengths of 512k, 1024k, and 2048k via PI and YaRN for comparison. As demonstrated in Table~\ref{tbl:secondsearch}, our non-uniform positional interpolation technique maintains a steady perplexity across different extension ratios. Conversely, when employing PI or YaRN, perplexity escalates rapidly as the context length increases.

\textbf{Recovery effectiveness at reduced context lengths.} To address the drop in performance associated with shorter context lengths, we fine-tune the RoPE parameters for LongRoPE-2048k by employing our search algorithm. In particular, we constrain the maximum scale factors permitted during the search, which serves to reduce the extent of interpolation at 4k and 8k context lengths. Table~\ref{tbl:shortperformance} presents a perplexity analysis of LongRoPE-LLaMA2-2048k on Proof-pile at these shorter lengths, together with the mean LLM benchmark accuracy. These findings indicate a marked enhancement in performance at reduced context lengths.

\textbf{Analysis on the two forms of non-uniformities}. To further investigate the effects of the two types of non-uniformities, we perform ablation studies assessing the contribution of each to model performance. Two experimental setups are considered: (i) extending LLaMA2-7B to shorter sequence lengths of 16k and 32k through various approaches—namely PI, optimizing only over the RoPE dimension, and jointly optimizing both non-uniformities; and (ii) expanding the context window of a 256k-length fine-tuned LLaMA2 to 2048k using analogous strategies. All models are assessed on perplexity scores in the absence of additional fine-tuning. Results in Table~\ref{tbl:searchdimension} indicate that introducing non-uniformity in the RoPE dimension offers notable perplexity improvements over the linear interpolation used in PI. Adjusting token position non-uniformity leads to performance gains for 16k and 32k context lengths but appears less effective for 2048k, likely because of the challenge posed by such an extensive sequence. Maintaining the original tokens without interpolation loses effectiveness under these circumstances, and addressing this limitation is deferred to future research.

\section{Related Works}

Beyond position interpolation strategies, this section reviews other relevant methodologies.

\textbf{Retrieval-based} techniques utilize external memory components to store historical context and employ retrieval mechanisms to access pertinent documents during inference~\cite{longllama,wang2023augmenting,borgeaud2022improving}. Such frameworks generally require significant adjustments to the underlying LLM models. By contrast, our approach introduces only minor modifications to the positional embeddings, making it a more lightweight solution. Additionally, our method extends to a broader set of long-context applications beyond retrieval, including long document summarization and few-shot learning.

\textbf{Attention-based context window extensions}. In addition to extending context through positional embedding interpolation, certain studies have expanded the input context window for LLMs by altering attention mechanisms within the original context window size~\cite{han2023lminfinite, streamingllm,parallelwindow}. Central to these approaches is the development of innovative attention masks to address the attention explosion problem introduced by additional positions. These approaches complement those based on positional interpolation.

\textbf{Fine-tuning based approaches} concentrate on strategies for adapting pre-trained LLMs to support longer contexts by adjusting position embeddings. Approaches such as Code LLaMA~\cite{codellama}, LLaMA2 Long~\cite{xiong2023effective}, and ScaledRoPE~\cite{liu2023scaling} employ significantly enlarged base values for RoPE and perform fine-tuning on sequences of the intended length. In contrast, our technique is adaptable across a range of target sequence lengths and is capable of handling inputs beyond 2M tokens. As fine-tuning for extended contexts (exceeding 128k tokens) is highly resource-intensive, methods like LongLoRA~\cite{longlora} and PoSE~\cite{zhu2023pose} have been introduced to alleviate this computational burden. Our approach remains complementary and independent from these resource-efficient fine-tuning solutions.

\section{Conclusion}

In this study, we introduce LongRoPE, a technique that dramatically increases the context window of LLMs to an unprecedented 2048k tokens, all while preserving their effectiveness within the original, shorter context range. Our approach leverages two types of non-uniformities present in RoPE positional embeddings through a computationally efficient evolutionary search process. This methodology yields two main advantages: it supplies strong initialization for subsequent fine-tuning and allows for an 8$\times$ expansion of the context window even in the absence of fine-tuning. Building upon these insights, we develop a progressive extension strategy whereby LLMs fine-tuned with a 256k context can be expanded to a 2048k context window without additional fine-tuning. Comprehensive empirical evaluations demonstrate the robustness and utility of LongRoPE. We anticipate that the LongRoPE-2048k models will unlock numerous new long-context use cases and encourage future research in this area.

\section*{Broader Impacts} The research described in this paper aims to further developments within the discipline of Machine Learning. While our contributions may have various implications for society, we do not identify any particular consequences that warrant special attention in this context.

	\balance

\end{document}